\section{Finite source estimates on the shifted charts}\label{me:analytic}
The exact moving finite calculus is imported from
\cite[Proposition 6.1, Theorem 7.1 and Proposition 8.1]{Meridian}:
these are respectively the full crossing polynomial and trace,
the original tree contraction, and its scalar ledger. In its fixed
Newton basis, with $d=N-1$,
\[
 e_k(z)=S_k^{-2}\prod_{j<k}(c_j-z),\quad
 S_k=\prod_{r=1}^k2\sin(\pi r/N),\quad
 c_j=2\cos((2j+1)\pi/N),
\]
the quotient is $P_\sigma=\mathsf T_N+2\cos2\pi\sigma$,
$\mathsf T_N(u+u^{-1})=u^N+u^{-N}$. The metric is
$\phi_\sigma(e_i e_j)=\mathbf1_{i+j=d}$ with
$\phi_\sigma=(-1)^dN^2[z^d]\operatorname{rem}_{P_\sigma}$.
An original physical edge, in lower-index convention, has equal
indices and coefficient $\gamma_{e,N}(s)H_N(s)/W_k(s)$, where
\begin{equation}\label{me:W}
 W_k(s)=\prod_{r=1}^k
       \left(1-\frac{\sin^2(\pi s/N)}{\sin^2(\pi r/N)}\right),
 \qquad\gamma_{e,N}\in\{1,q^{s^2/2},q^{-s^2/2}\}.
\end{equation}
The complete Hermite norm uses all distinct node values and the
first derivative at each double node. The odd node $-2$ is simple.
Polynomial equivalence with interval and Chebyshev norms is
\cite[Lemma 2.1]{Bounds}; quotient products of fixed valence are
first reduced, so the assertion does not identify the full matrix
algebra with a commutative polynomial algebra.

\begin{lemma}[Translated compact sets]\label{me:translated}
For every source class in Section~\ref{me:intro}, the local selected
source, its nonzero leading coefficient, and all fixed exponential
error margins hold on a fixed compact $K\subset\{\Im s>0\}$.
Only bounded or polynomial prefactors and starting levels change.
The same is true on a finite union of translates $K+j$.
\end{lemma}
\begin{proof}
For strict Conway sources this is part of
\cite[Theorems 1.1--1.2]{MovingSources}. We give the finite argument
also needed for the all-circle classes. Euler's sine product gives
\begin{equation}\label{me:subset}
 W_k(s)=\prod_{\substack{1\le r\le k\\a\in\mathbb Z}}
                  \left(1-\frac{s^2}{(r+aN)^2}\right).
\end{equation}
The nonzero integers are distinct. On $K\cup(-K)$ the finitely
many small factors are bounded above and away from zero, and the
logarithms of all remaining factors are bounded by $C_K/m^2$.
Products of the minima with one and maxima with one give bounds
for every subset, uniformly in $N,k$. Thus $W_k$, its inverse,
$H_N$ and its inverse are uniformly bounded. Positivity in the
all-circle estimates is used only for the complete Newton factors;
these bounded multipliers do not change their exponential budgets.

For clarity the exact fixed-power interpolation is
\begin{equation}\label{me:crossing}
 r_{N,\sigma}^{[n]}(z)=
 -\frac1{N(U-U^{-1})}\sum_{j=0}^{N-1}
 \lambda_j^n(t_j-t_j^{-1})\frac{P_\sigma(z)}{z-c_j(\sigma)},
\end{equation}
where $U=e^{2\pi\ii\sigma}$,
$t_j=q^{-2j-1+2\sigma}$, $c_j(\sigma)=t_j+t_j^{-1}$ and
$\lambda_j=(-1)^jq^{j(j+1)-2\sigma j}$.
The full power includes $\kappa_\sigma^n$ with
$\kappa_\sigma=q^{(d+\sigma)^2/2}$. The denominator is bounded
away from zero on the nonreal compact. Nodes are in a Chebyshev
ellipse of angular height $O_K(1/N)$. The divided difference and
its derivative have polynomial bounds there, while the fixed
positive or negative powers of the eigenvalues and carrier have
bounded modulus. This proves the same full-polynomial bounds as
on a disk. Exact Chebyshev folding costs a bounded factor per
fold; the number of folds is fixed by the vertex valence.

At a constant opposite-color vertex keep a separate formal
$\sigma$ and use its whole trace on a parameter circle avoiding
collisions. Its confluent normals are removable by the finite
crossing theorem. The maximum principle for the whole trace
minus its finite approximation brings this formal parameter to
zero, while the edge parameter $s$ stays independent. This is not
an identification of constant and moving orientation types.

The actual trace grid is shifted by $2\pi\sigma/N$. Its exact
pole kernel has vertical-side value
$\ii N e^{-Nv}/(1+e^{-Nv})$ after that translation. Thus the
endpoint action loss survives the $O_K(1/N)$ displacement. A
translated chart $s+j$ may wrap finitely many grid points across
$\pm\pi$; retain those entries and their chirps. Their endpoint
loss persists, but the chirp itself is not periodic. Paired charge
coefficients are
\begin{equation}\label{me:charge}
 (-1)^{d(\ell-1)}U_{\ell-1}(\cos2\pi\sigma),\qquad
 |U_{\ell-1}(\cos2\pi\sigma)|
       \le\ell e^{2\pi|\Im\sigma|(\ell-1)}.
\end{equation}
The added fixed exponential in $\ell$ is absorbed by the original
$e^{-N\ell y}$ tail at every fixed positive contour height $y$.
The first coefficient stays exactly one. Real/strip estimates
in \cite[Sections 2--6]{Summed} and the conditional estimates in
\cite[Sections 3--7]{Directed} use precisely these bounded
multipliers, polynomial finite norms and geometric tails. All
exponential margins therefore persist. This verifies their source
hypotheses on $K$ directly; it is not analytic continuation of
an asymptotic estimate through resonant real parameters.
\end{proof}

\begin{figure}[t]
\centering\begin{tikzpicture}[x=1.55cm,y=1.3cm,>=Latex,every node/.style={font=\small}]
\draw[->](-.6,0)--(6.2,0) node[right]{$\operatorname{Re}s$};
\draw[->](-.5,-.2)--(-.5,1.8) node[above]{$\operatorname{Im}s$};
\foreach \j in {0,1,2,3,4,5}{\fill(\j,0)circle(1.1pt);}
\foreach \j/\name in {0/D,1/{D+1},2/{D+2},5/{D+r_P}}{
 \draw[fill=blue!8] (\j,1)ellipse(.27 and .32);
 \node[above]at(\j,1.35){$\name$};
}
\node at(3.5,1){$\cdots$};
\draw[->,thick](.3,.95)--(.7,.95);
\node[below]at(.5,.48){$s\mapsto s+1$};
\node[below]at(0,0){$0$};\node[below]at(1,0){$1$};
\end{tikzpicture}

\caption{A fixed disk $D$ above the real parameter axis and its
translates $D+j$ provide the arguments in an order-$r_P$ recurrence.
The real dots are integer resonances; no contour crosses them.
The physical meridians are $M=q_N^s$, so every fixed disk still
shrinks to $M=1$ as $N$ grows.}
\label{me:chart-figure}
\end{figure}

\section{Canonical local integrals to every order}\label{me:local-orders}
The canonical factorial is
\begin{equation}\label{me:factorial}
 G_N(z)=\left(\frac{2\pi}{N}\right)^z\Gamma(z+1)
 \prod_{a\ge1}\frac{\Gamma(Na+z+1)}
                      {\Gamma(Na-z)(Na)^{2z+1}}.
\end{equation}
We use the complete proof in \cite[Lemma 2.1]{Asymptotics},
including its fixed normalization, sine shift and reflection:
\[
 G_N(k)=S_k,\qquad G_N(z)G_N(N-1-z)=N,
 \qquad G_N(s)G_N(-s)=H_N(s)^{-1}.
\]
On every smaller complex bulk compact it supplies, with all fixed
derivatives and every requested finite order,
\begin{equation}\label{me:factorial-expansion}
 \log G_N(Nx)=NJ(x)+\tfrac12\log(2N\sin\pi x)
           +\sum_{j=1}^{L}N^{-j}b_j(x)+O(N^{-L-1}).
\end{equation}
This is an asymptotic expansion, with no convergence assertion.
The summable Gamma remainders in that lemma justify the infinite
product before differentiation.

The exact local circle factor and the physical edge insertion are
\begin{align}
 \mathcal E_N(x,\tau)&=
 \frac{G_N(d/2+u)G_N(d/2-u)}
 {G_N(d/2-Nx+u)G_N(d/2-Nx-u)G_N(Nx)^2},
 \quad u=\frac{N\tau}{2\pi},\label{me:E}\\
 \mathcal I_N(x,s)&=\frac{G_N(Nx)^2}{G_N(Nx-s)G_N(Nx+s)}.
                                                        \label{me:I}
\end{align}
The sine shift proves their equality with
$e_k(-2\cos\tau)$ and $H_N/W_k$ at $x=k/N$.
In particular the reflection factor in \eqref{me:factorial} puts
$H_N$ in the numerator of the edge weight. Hidden complete gates
are ratios of the same factorials. Every argument, including
$d/2=N/2-1/2$ and every $d-k$ offset, remains exact.

\begin{lemma}[Uniform local expansion]\label{me:local-expansion}
On the selected primitive wire contours and the selected local
angular contours, the canonical integrals have expansions to every
relative inverse-level order. Their coefficients and remainders
are holomorphic on the stated smaller parameter compacts. This
applies to a conditional edge integral with its parent angle still
free, and to a joint strict Conway or summed-circle integral.
\end{lemma}
\begin{proof}
All factorial arguments lie in fixed bulk complex tubes on these
contours. For strict sources this is proved in
\cite[Section 6]{MovingSources}; hidden unturned indices need not
be below half. The exceptional half physical edge is bulk and
causes no pole. For all-circle sources the edge saddle is in the
fixed tube about $1/5<x<3/10$, as proved in
\cite[Section 3]{Summed} and \cite[Section 3]{Directed}.
Substitute \eqref{me:factorial-expansion} into the exact ratios.
Powers of $N$, integer offsets and square-root marks give the
original leading amplitude and complete inverse-power expansions.

The finite lattice is handled before asymptotic integration.
Enclose each primitive coordinate interval by half-grid rectangles
and use $\pi N\cot(\pi Nz)$. On horizontal sides the difference
from its constant value is $O(e^{-2\pi Nh})$; on vertical sides
it is bounded. Choose the height $h$ before the parameter tube,
with the phase gain $C h^2$ less than both a fixed fraction of the
boundary loss and $\pi h$. Thus nonconstant modes and boundary
currents have a fixed exponential loss. The constant terms give
exactly $N^K$ times the local integral, up to those errors. The
primitive lattice and all offsets are the ones in the imported
finite source, so no sublattice multiplicity is introduced.

Translate the local contour to its selected critical point. Its
real quadratic part is uniformly negative on compact parameter
sets; its inverse Hessian and leading amplitude are bounded and
nonzero there. To obtain relative order $L$, Taylor expand phase
and amplitude to a sufficiently high finite order, set
$X-X_*=N^{-1/2}Y$, and split at $|X-X_*|=N^{-\epsilon}$,
with $0<\epsilon<1/6$ chosen first. Outside this shrinking region
the quadratic bound is $e^{-cN^{1-2\epsilon}}$. Inside it expand
the exponential of the cubic and higher terms with its integral
remainder. Taking enough Taylor terms makes that remainder at most
$N^{-L-1}$ times an integrable polynomial times
$e^{-c|Y|^2/2}$. The finitely many Gaussian moments give integer
relative powers: odd total moments vanish. Fixed derivatives are
obtained by starting with more terms and using a smaller tube.

The absolute Gaussian integral is bounded by a constant times the
modulus of its nonzero leading complex integral. Thus a uniform
holomorphic relative remainder stays of the same order after
integration. The bound may depend on fixed degrees and powers,
which is permitted. This proves the claim directly, rather than
promoting a first-order estimate to all orders without a remainder
argument. Conditional angular and edge integrations use the same
proof with the unevaluated parent coordinate as parameter.
\end{proof}

For the strict Conway and summed-circle classes the fixed global
source and charge gaps are already below the same joint first
charge. Lemmas~\ref{me:translated}--\ref{me:local-expansion}
therefore give their entire source to every relative order. The
source error is kept additive until that comparison; it is never
compared with a product of separately damped local contributions.
The directed case requires one additional normalized induction,
which we now give.
